\documentclass{ctexart}
\usepackage[dvipsnames]{xcolor}
\usepackage{tikz}
\usepackage{pgfplots}
\usetikzlibrary{arrows.meta,graphs,graphdrawing}
\usegdlibrary{layered}
\pgfplotsset{compat=1.18}
\usepackage{amssymb}
\usepackage{siunitx}

\newcommand\bbQ{\mathbb{Q}}

\begin{document}
\[
\begin{tikzpicture}
\draw (-3,0) -- (3,0);
\draw (0,-3) -- (0,3);
\end{tikzpicture}
\]
\[
\begin{tikzpicture}
\draw[->] (-3,0) -- (3,0);
\draw[->] (0,-3) -- (0,3);
\end{tikzpicture}
\]
\[
  \begin{tikzpicture} \draw [->] (-3,0) -- (3,0);
    \draw [->] (0,-3) -- (0,3);
    \draw [smooth, samples=100, domain=-1.5:1.5] plot(\x, {(\x)^2)});
    \draw [smooth, samples=100, domain=-1.5:1.5] plot(\x, {3*(\x)^3-4*(\x)});
    \end{tikzpicture}
\]
\[
  \begin{tikzpicture}
    \draw [->] (-3,0) -- (3,0);
    \draw [->] (0,-3) -- (0,3);
    \draw [smooth, samples=100, domain=-1.6:1.6] plot(\x, {(\x)^2)}); \draw [green, smooth, samples=5, domain=-1.6:1.6] plot(\x, {(\x)^2)});
    \draw [red, samples=5, domain=-1.6:1.6] plot(\x, {(\x)^2)}); \end{tikzpicture}
\]
\[
  \begin{tikzpicture} \draw [->] (-3,0) -- (3,0);
    \draw [->] (0,-3) -- (0,3);
    \draw [smooth, samples=100, domain=-1.6:1.6] plot(\x, {(\x)^2)});
    \node at (2.0,1.8) {$y=x^2$}; \end{tikzpicture}
\]
\[
% a picture of the parabola y=x^2
\begin{tikzpicture}
  % the x-axis, with an arrow pointing right
  \draw [->] (-3,0) -- (3,0);

  % the y-axis, with an arrow pointing up
  \draw [->] (0,-3) -- (0,3);

  % the parabola y=x^2, from -1.6 to 1.6,
  % using sample size = 100 points 
  % and joining points smoothly 
  \draw [smooth, samples=100, domain=-1.6:1.6] plot(\x, {(\x)^2)});

  % the label y=x^2 
  \node at (2.0,1.8) {$y=x^2$};
  \end{tikzpicture}
\]
\[
\begin{tikzpicture}
    % the x-axis, with an arrow pointing right
    \draw [->] (-3,0) -- (3,0);

    % the y-axis, with an arrow pointing up
    \draw [->] (0,-3) -- (0,3);
  \draw (-1,1) -- (1,1);
  \draw (1,1)  -- (1,-1);
  \draw (-1,-1) -- (1,-1);
  \draw (-1,-1) -- (-1,1);
\end{tikzpicture}
\]
\[
\begin{tikzpicture}
    % the x-axis, with an arrow pointing right
    \draw [->] (-3,0) -- (3,0);

    % the y-axis, with an arrow pointing up
    \draw [->] (0,-3) -- (0,3);

    \draw (0,1) -- (1,0);
    \draw (1,0) -- (0,-1);
    \draw (0,-1) -- (-1,0);
    \draw (-1,0) -- (0,1);

\end{tikzpicture}
\]

\begin{center}
  \begin{tikzpicture}

      \filldraw[fill=yellow, draw=black] (0,1) -- (1,0) -- (0,-1) -- (-1,0) -- cycle;
    % the x-axis, with an arrow pointing right
    \draw [->] (-3,0) -- (3,0);

    % the y-axis, with an arrow pointing up
    \draw [->] (0,-3) -- (0,3);
  \end{tikzpicture}
\end{center}

\[
  \begin{tikzpicture} \draw [->] (-3,0) -- (3,0);
    \draw [->] (0,-7) -- (0,3);
    \draw [smooth, samples=100, domain=-1.6:1.6] plot(\x, {-(\x)^3)+2*(\x)-5});
    \node at (3.0,-5) {$y=-x^3+2x-5$}; 
  \end{tikzpicture}
\]

\[
  \begin{tikzpicture}
    \draw [->] (-3,0) -- (3,0);
    \draw [->] (0,-3) -- (0,3);
    \draw [red, smooth, samples=100, domain=-1.6:1.6] plot(\x, {(\x)^2)-2}); 
    \draw [blue, smooth, samples=100, domain=-1.6:1.6] plot(\x, {-(\x)^2+1)});
    \node at (1,-2.1) {$y=x^2-2$}; 
    \node at (1.4,1) {$y=-x^2+1$}; 
  \end{tikzpicture}
\]
\[
  \begin{tikzpicture} 
    \draw [help lines] (-1,-4) grid (5,4);
  \end{tikzpicture}
\]
\[
\begin{tikzpicture}
% help lines (comment out when done) 

\draw [help lines] (-1,-4) grid (5,4);
% edges from vertex a
\draw (0,2) -- (4,3);
\draw (0,2) -- (4,1); \draw (0,2) -- (4,-3);
% edges from vertex b
\draw (0,0) -- (4,1); \draw (0,0) -- (4,-1);
% edges from vertex c \draw (0,-2) -- (4,3);
\draw (0,-2) -- (4,1);
\end{tikzpicture}
\]

\[
  \begin{tikzpicture}
    % help lines (comment out when done) 
    \draw [help lines] (-1,-4) grid (5,4);
    % edges from vertex a
    \draw (0,2) -- (4,3);
    \draw (0,2) -- (4,1); \draw (0,2) -- (4,-3);
    % edges from vertex b
    \draw (0,0) -- (4,1); \draw (0,0) -- (4,-1);
    % edges from vertex c
    \draw (0,-2) -- (4,3);
    \draw (0,-2) -- (4,1);
    % dot at vertex a
    \node [circle, fill, inner sep=2pt] at (0,2) {};
    \end{tikzpicture}
\]

\[
  \begin{tikzpicture}
    % vertex style to be used at vertex nodes
    \tikzstyle{vertex}=[circle, fill, inner sep=2pt]
    % help lines (comment out when done) 
    \draw [help lines] (-1,-4) grid (5,4);
    % edges from vertex a
    \draw (0,2) -- (4,3); \draw (0,2) -- (4,1); \draw (0,2) -- (4,-3);
    % edges from vertex b
    \draw (0,0) -- (4,1); \draw (0,0) -- (4,-1);
    % edges from vertex c
    \draw (0,-2) -- (4,3); \draw (0,-2) -- (4,1);
    % dot at vertex a 
    \node [vertex] at (0,2) {};
    \end{tikzpicture}
\]

\[
  \begin{tikzpicture}
    % vertex style to be used at vertex nodes
    \tikzstyle{vertex}=[circle, fill, inner sep=2pt]
    % help lines (comment out when done) 
    \draw [help lines] (-1,-4) grid (5,4);
    % edges from vertex a
    \draw (0,2) -- (4,3); \draw (0,2) -- (4,1); \draw (0,2) -- (4,-3);
    % edges from vertex b
    \draw (0,0) -- (4,1); \draw (0,0) -- (4,-1);
    % edges from vertex c 
    \draw (0,-2) -- (4,3);
    \draw (0,-2) -- (4,1);
    % dot at vertex a 
    \node [vertex] at (0,2) {};
    % dot at vertex b 
    \node [vertex] at (0,0) {};
    % dot at vertex c 
    \node [vertex] at (0,-2) {};
    % dot at vertex d 
    \node [vertex] at (4,3) {};
    % dot at vertex e 
    \node [vertex] at (4,1) {};
    % dot at vertex f 
    \node [vertex] at (4,-1) {};
    % dot at vertex g 
    \node [vertex] at (4,-3) {};
\end{tikzpicture}
\]

\[
  \begin{tikzpicture}
    % vertex style to be used at vertex nodes 
    \tikzstyle{vertex}=[circle, fill, inner sep=2pt]
    % help lines (comment out when done)
    \draw [help lines] (-1,-4) grid (5,4);
    % edges from vertex a
    \draw (0,2) -- (4,3);
    \draw (0,2) -- (4,1); \draw (0,2) -- (4,-3);
    % edges from vertex b
    \draw (0,0) -- (4,1); \draw (0,0) -- (4,-1);
    % edges from vertex c 
    \draw (0,-2) -- (4,3);
    \draw (0,-2) -- (4,1);
    % dot at vertex a 
    \node [vertex] at (0,2) {};
    % label at vertex x 
    \node [left] at (0,2) {$a$};
% dot at vertex b 
\node [vertex] at (0,0) {};
% dot at vertex c 
\node [vertex] at (0,-2) {};
% dot at vertex d 
\node [vertex] at (4,3) {};
% dot at vertex e
\node [vertex] at (4,1) {};
% dot at vertex f 
\node [vertex] at (4,-1) {};
% dot at vertex g 
\node [vertex] at (4,-3) {};
\end{tikzpicture}
\]

\[
  \begin{tikzpicture}
    % vertex style to be used at vertex nodes
    \tikzstyle{vertex}=[circle, fill, inner sep=2pt]
    % help lines (comment out when done)
    \draw [help lines] (-1,-4) grid (5,4);
    % edges from vertex a
    \draw (0,2) -- (4,3); \draw (0,2) -- (4,1);
    \draw (0,2) -- (4,-3);
% edges from vertex b
\draw (0,0) -- (4,1); 
\draw (0,0) -- (4,-1);
% edges from vertex c 
\draw (0,-2) -- (4,3);
\draw (0,-2) -- (4,1);
% dot at vertex a 
\node [vertex] at (0,2) {};
% label at vertex a 
\node [left] at (0,2) {$a$};
% dot at vertex b 
\node [vertex] at (0,0) {};
% label at vertex b 
\node [left] at (0,0) {$b$};
% dot at vertex c 
\node [vertex] at (0,-2) {};
% label at vertex a 
\node [left] at (0,-2) {$c$};
% dot at vertex d
\node [vertex] at (4,3) {}; % label at vertex d 
\node [right] at (4,3) {$d$};
% dot at vertex e 
\node [vertex] at (4,1) {};
% label at vertex e 
\node [right] at (4,1) {$e$};
% dot at vertex f 
\node [vertex] at (4,-1) {};
% label at vertex f
\node [right] at (4,-1) {$f$};
% dot at vertex g 
\node [vertex] at (4,-3) {}; 
% label at vertex g 
\node [right] at (4,-3) {$g$};
\end{tikzpicture}
\]

\[
  \begin{tikzpicture}
    % vertex style to be used at vertex nodes
    \tikzstyle{vertex}=[circle, fill, inner sep=2pt]
    % help lines (comment out when done) 
    % \draw [help lines] (-1,-4) grid (5,4);
    % edges from vertex a
    \draw (0,2) -- (4,3); \draw (0,2) -- (4,1); \draw (0,2) -- (4,-3);
    % edges from vertex b
    \draw (0,0) -- (4,1); \draw (0,0) -- (4,-1);
    % edges from vertex c 
    \draw (0,-2) -- (4,3);
    \draw (0,-2) -- (4,1);
    % dot at vertex a 
    \node [vertex] at (0,2) {};
    % label at vertex a \node [left] at (0,2) {$a$};
    % dot at vertex b 
    \node [vertex] at (0,0) {};
    % label at vertex b 
    \node [left] at (0,0) {$b$};
    % dot at vertex c
    \node [vertex] at (0,-2) {};
    % label at vertex a 
    \node [left] at (0,-2) {$c$};
    % dot at vertex d
\node [vertex] at (4,3) {};
% label at vertex d 
\node [right] at (4,3) {$d$};
% dot at vertex e
\node [vertex] at (4,1) {};
% label at vertex e 
\node [right] at (4,1) {$e$};
% dot at vertex f 
\node [vertex] at (4,-1) {};
% label at vertex f
\node [right] at (4,-1) {$f$};
% dot at vertex g 
\node [vertex] at (4,-3) {}; 
% label at vertex g 
\node [right] at (4,-3) {$g$};
\end{tikzpicture}
\]


\begin{tikzpicture}
  % vertex style to be used at vertex nodes
  \tikzstyle{vertex}=[circle, fill, inner sep=2pt]
  % dot at vertex a
\node [vertex] at (0,2) {};
% label at vertex a, offset = 0 
\node [left=0pt] at (0,2) {$a$};
\end{tikzpicture}

\begin{tikzpicture}
% vertex style to be used at vertex nodes 
\tikzstyle{vertex}=[circle, fill, inner sep=2pt]
% dot at vertex a 
\node [vertex] at (0,2) {};
% label at vertex a, offset = 1 
\node [left=1pt] at (0,2) {$a$};
\end{tikzpicture}

\begin{tikzpicture}
% vertex style to be used at vertex nodes 
\tikzstyle{vertex}=[circle, fill, inner sep=2pt]
% dot at vertex a 
\node [vertex] at (0,2) {}; 
% label at vertex a, offset = 2
\node [left=2pt] at (0,2) {$a$};
\end{tikzpicture}

\begin{tikzpicture}
% vertex style to be used at vertex nodes 
\tikzstyle{vertex}=[circle, fill, inner sep=2pt]
% dot at vertex a 
\node [vertex] at (0,2) {}; 
% label at vertex a, offset = 3
\node [left=3pt] at (0,2) {$a$};
\end{tikzpicture}

\begin{tikzpicture}
% vertex style to be used at vertex nodes
\tikzstyle{vertex}=[circle, fill, inner sep=2pt]
% dot at vertex a 
\node [vertex] at (0,2) {};
% label at vertex a, offset = 4
\node [left=4pt] at (0,2) {$a$};
\end{tikzpicture}

\begin{tikzpicture}
  % vertex style to be used at vertex nodes
  \tikzstyle{vertex}=[circle, fill, inner sep=2pt]
  % help lines (comment out when done) 
  % \draw [help lines] (-1,-4) grid (5,4);
  % edges from vertex a
  \draw (0,2) -- (4,3);
  \draw (0,2) -- (4,1); \draw (0,2) -- (4,-3);
  % edges from vertex b
  \draw (0,0) -- (4,1); \draw (0,0) -- (4,-1);
  % edges from vertex c
  \draw (0,-2) -- (4,3);
  \draw (0,-2) -- (4,1);
  % dot at vertex a 
  \node [vertex] at (0,2) {};
  % label at vertex a 
  \node [left=4pt] at (0,2) {$a$};
  % dot at vertex b 
  \node [vertex] at (0,0) {}; 
  % label at vertex b 
  \node [left=4pt] at (0,0) {$b$};
  % dot at vertex c
  \node [vertex] at (0,-2) {}; 
  % label at vertex a 
  \node [left=4pt] at (0,-2) {$c$};
  % dot at vertex d
  \node [vertex] at (4,3) {};
  % label at vertex d 
  \node [right=4pt] at (4,3) {$d$};
  % dot at vertex e 
  \node [vertex] at (4,1) {};
  % label at vertex e 
  \node [right=4pt] at (4,1) {$e$};
% dot at vertex f 
\node [vertex] at (4,-1) {};
% label at vertex f 
\node [right=4pt] at (4,-1) {$f$};
% dot at vertex g 
\node [vertex] at (4,-3) {}; 
% label at vertex g 
\node [right=4pt] at (4,-3) {$g$};
\end{tikzpicture}


\begin{tikzpicture}
  % vertex style to be used at vertex nodes
  \tikzstyle{vertex}=[circle, fill, inner sep=2pt]
  % help lines (comment out when done) 
  \draw [help lines] (-2,-2) grid (2,2);

  \node [vertex] at (-2,-2) {};
  \node [vertex] at (-2,-1) {};
  \node [vertex] at (-2,0) {};
  \node [vertex] at (-2,1) {};
  \node [vertex] at (-2,2) {};

  \node [vertex] at (-1,-2) {};
  \node [vertex] at (-1,-1) {};
  \node [vertex] at (-1,0) {};
  \node [vertex] at (-1,1) {};
  \node [vertex] at (-1,2) {};

  \node [vertex] at (0,-2) {};
  \node [vertex] at (0,-1) {};
  \node [circle,fill, red, inner sep=4pt] at (0,0) {};
  \node [vertex] at (0,1) {};
  \node [vertex] at (0,2) {};

  \node [vertex] at (1,-2) {};
  \node [vertex] at (1,-1) {};
  \node [vertex] at (1,0) {};
  \node [vertex] at (1,1) {};
  \node [vertex] at (1,2) {};

  \node [vertex] at (2,-2) {};
  \node [vertex] at (2,-1) {};
  \node [vertex] at (2,0) {};
  \node [vertex] at (2,1) {};
  \node [vertex] at (2,2) {};
\end{tikzpicture}

\begin{tikzpicture}
  % vertex style to be used at vertex nodes
  \tikzstyle{vertex}=[circle, fill, inner sep=2pt]
  % help lines (comment out when done) 
  % \draw [help lines] (-1,-4) grid (5,4);
  % edges from vertex a
  \draw (0,2) -- (4,3);
  \draw (0,2) -- (4,1); \draw (0,2) -- (4,-3);
  % edges from vertex b
  %\draw (0,0) -- (4,1); 
  \draw (0,0) -- (4,-1);
  % edges from vertex c
  \draw (0,-2) -- (4,3);
  \draw (0,-2) -- (4,1);

  \draw (0,2) -- (4,-1);
  % dot at vertex a 
  \node [vertex] at (0,2) {};
  % label at vertex a 
  \node [left=4pt] at (0,2) {$a$};
  % dot at vertex b 
  \node [vertex] at (0,0) {}; 
  % label at vertex b 
  \node [left=4pt] at (0,0) {$b$};
  % dot at vertex c
  \node [vertex] at (0,-2) {}; 
  % label at vertex a 
  \node [left=4pt] at (0,-2) {$c$};
  % dot at vertex d
  \node [vertex] at (4,3) {};
  % label at vertex d 
  \node [right=4pt] at (4,3) {$d$};
  % dot at vertex e 
  \node [vertex] at (4,1) {};
  % label at vertex e 
  \node [right=4pt] at (4,1) {$e$};
% dot at vertex f 
\node [vertex] at (4,-1) {};
% label at vertex f 
\node [right=4pt] at (4,-1) {$f$};
% dot at vertex g 
\node [vertex] at (4,-3) {}; 
% label at vertex g 
\node [right=4pt] at (4,-3) {$g$};
\end{tikzpicture}

\begin{tikzpicture}
  % vertex style to be used at vertex nodes
  \tikzstyle{vertex}=[circle, fill, inner sep=2pt]
  % help lines (comment out when done) 
  % \draw [help lines] (-1,-5) grid (5,4);
  % edges from vertex a
  \draw (0,2) -- (4,3);
  \draw (0,2) -- (4,1); 
  \draw (0,2) -- (4,-3);
  % edges from vertex b
  \draw (0,0) -- (4,1); 
  \draw (0,0) -- (4,-1); 
  \draw (0,0) -- (4,-5);
  % edges from vertex c
  \draw (0,-2) -- (4,3);
  \draw (0,-2) -- (4,1);
  \draw (0,-2) -- (4,-5);
  % dot at vertex h
  \node [vertex] at (4,-5) {};
  % label at vertex h
  \node [right=4pt] at (4,-5) {$h$};
  % dot at vertex a 
  \node [vertex] at (0,2) {};
  % label at vertex a 
  \node [left=4pt] at (0,2) {$a$};
  % dot at vertex b 
  \node [vertex] at (0,0) {}; 
  % label at vertex b 
  \node [left=4pt] at (0,0) {$b$};
  % dot at vertex c
  \node [vertex] at (0,-2) {}; 
  % label at vertex a 
  \node [left=4pt] at (0,-2) {$c$};
  % dot at vertex d
  \node [vertex] at (4,3) {};
  % label at vertex d 
  \node [right=4pt] at (4,3) {$d$};
  % dot at vertex e 
  \node [vertex] at (4,1) {};
  % label at vertex e 
  \node [right=4pt] at (4,1) {$e$};
% dot at vertex f 
\node [vertex] at (4,-1) {};
% label at vertex f 
\node [right=4pt] at (4,-1) {$f$};
% dot at vertex g 
\node [vertex] at (4,-3) {}; 
% label at vertex g 
\node [right=4pt] at (4,-3) {$g$};
\end{tikzpicture}

\begin{tikzpicture}
  % vertex style to be used at vertex nodes
  \tikzstyle{vertex2}=[circle, fill, green, inner sep=1pt]
  % help lines (comment out when done) 
  % \draw [help lines] (-1,-4) grid (5,4);
  % edges from vertex a
  \draw (0,2) -- (4,3);
  \draw (0,2) -- (4,1); \draw (0,2) -- (4,-3);
  % edges from vertex b
  \draw (0,0) -- (4,1); \draw (0,0) -- (4,-1);
  % edges from vertex c
  \draw (0,-2) -- (4,3);
  \draw (0,-2) -- (4,1);
  % dot at vertex a 
  \node [vertex2] at (0,2) {};
  % label at vertex a 
  \node [left=4pt] at (0,2) {$a$};
  % dot at vertex b 
  \node [vertex2] at (0,0) {}; 
  % label at vertex b 
  \node [left=4pt] at (0,0) {$b$};
  % dot at vertex c
  \node [vertex2] at (0,-2) {}; 
  % label at vertex a 
  \node [left=4pt] at (0,-2) {$c$};
  % dot at vertex d
  \node [vertex2] at (4,3) {};
  % label at vertex d 
  \node [right=4pt] at (4,3) {$d$};
  % dot at vertex e 
  \node [vertex2] at (4,1) {};
  % label at vertex e 
  \node [right=4pt] at (4,1) {$e$};
% dot at vertex f 
\node [vertex2] at (4,-1) {};
% label at vertex f 
\node [right=4pt] at (4,-1) {$f$};
% dot at vertex g 
\node [vertex2] at (4,-3) {}; 
% label at vertex g 
\node [right=4pt] at (4,-3) {$g$};
\end{tikzpicture}

\[
  \begin{tikzpicture}
    % vertex style to be used at vertex nodes
    \tikzstyle{vertex}=[circle, fill, inner sep=2pt]
    % help lines (comment out when done) 
    % \draw [help lines] (-4,-3) grid (4,3);

    \node [vertex] at (-4,3) {};
    % label at vertex a 
    \node [above=4pt] at (-4,3) {$a$};

    \node [vertex] at (0,3) {};
    % label at vertex b
    \node [above=4pt] at (0,3) {$b$};

    \node [vertex] at (4,3) {};
    % label at vertex c
    \node [above=4pt] at (4,3) {$c$};


    \node [vertex] at (-4,-3) {};
    % label at vertex d
    \node [below=4pt] at (-4,-3) {$d$};

    \node [vertex] at (0,-3) {};
    % label at vertex e
    \node [below=4pt] at (0,-3) {$e$};

    \node [vertex] at (4,-3) {};
    % label at vertex f
    \node [below=4pt] at (4,-3) {$f$};

    % edges from vertex a
    \draw (-4,3) -- (-4,-3);
    \draw (-4,3) -- (0,-3); 
    \draw (-4,3) -- (4,-3);
    % edges from vertex b
    \draw (0,3) -- (-4,-3);
    \draw (0,3) -- (0,-3); 
    \draw (0,3) -- (4,-3);
    % edges from vertex c
    \draw (4,3) -- (-4,-3);
    \draw (4,3) -- (0,-3); 
    \draw (4,3) -- (4,-3);

  \end{tikzpicture}
\]  


\begin{tikzpicture} \graph [layered layout]{ 
  a -- {b, c, d}; 
  {b, c, d} -- e; 
  };
\end{tikzpicture}

\begin{tikzpicture}
  \graph [layered layout] {
  "$\mathbb{Q}(\sqrt{2}, \sqrt{3})$" --
  {"$\mathbb{Q}(\sqrt{2})$", "$\mathbb{Q}(\sqrt{6})$", 
  "$\mathbb{Q}(\sqrt{3})$"}; 
  {"$\mathbb{Q}(\sqrt{2})$", "$\mathbb{Q}(\sqrt{6})$", 
  "$\mathbb{Q}(\sqrt{3})$"}
  -- "$\mathbb{Q}$";
  }; 
\end{tikzpicture}

\begin{tikzpicture} \graph [layered layout]{
  a -- {b, c, d};
  e -- {b, c, d}; }; 
\end{tikzpicture}

\begin{tikzpicture}
    \graph [layered layout]{
    a -- {b, c, d}; e -- {b, c, d}; { [same layer] a };
    { [same layer] b, c, d };
    { [same layer] e }; }; 
\end{tikzpicture}

\begin{tikzpicture}
      \graph [layered layout]{ a -- {b, "the quick red fox jumped over the lazy brown dog", d};
      {b, "the quick red fox jumped over the lazy brown dog", d} -- e; }; 
\end{tikzpicture}

\begin{tikzpicture}
  % nodes on the same level are spaced farther apart 
  \graph [layered layout, sibling distance=2cm] {
  a -- {b, c, d};
  {b, c, d} -- e; }; 
\end{tikzpicture}

% layers are spaced farther apart 
\begin{tikzpicture} \graph [layered layout, level distance=2cm] {
  a -- {b, c, d};
  {b, c, d} -- e; }; 
\end{tikzpicture}

\begin{tikzpicture}
  % graph is oriented vertically
  % (levels are vertical)
  \graph [layered layout, vertical = b to d] {
  a -- {b, c, d}; {b, c, d} -- e; }; 
\end{tikzpicture}

\begin{tikzpicture}
  % graph is oriented vertically
  % (levels are vertical)
  \graph [layered layout, vertical' = b to d] {
  a -- {b, c, d};
  {b, c, d} -- e;
  }; 
\end{tikzpicture}

\begin{tikzpicture}
  % node e is nudged a bit off-center
  \graph [layered layout]{ 
    a -- {b, c, d}; {b, c, d} -- e[nudge=(left:3mm)];
  }; 
\end{tikzpicture}

  \begin{tikzpicture} 
    \graph [layered layout, sibling distance=2cm,level distance=0.5cm] {
    a -- {b, c, d};
    {b, c, d} -- e; }; 
  \end{tikzpicture}

  \begin{tikzpicture} 
    \graph [layered layout] {
      {g,h,i} -- e;
      {e,f} -- c;
      {b,c,d} -- a;
    };
  \end{tikzpicture}


  \begin{tikzpicture} 
    \graph [layered layout] {
      "$D_8$" -- {"$\langle s,r^2\rangle$","$\langle r\rangle$","$\langle rs,r^2\rangle$"};
      "$\langle s,r^2\rangle$" -- {"$\langle s\rangle$","$\langle r^2s\rangle$","$\langle r^2\rangle$"};
      "$\langle r\rangle$" -- "$\langle r^2\rangle$";
      "$\langle rs,r^2\rangle$" -- {"$\langle r^2\rangle$","$\langle rs\rangle$","$\langle r^3s\rangle$"};
      {"$\langle s\rangle$","$\langle r^2s\rangle$","$\langle r^2\rangle$","$\langle rs\rangle$","$\langle r^3s\rangle$"} -- "$1$";
    };
  \end{tikzpicture}
  \[
    \begin{tikzpicture} 
      \graph [layered layout,level distance=2cm] {
        "$\bbQ(i,\sqrt[8]{2})$" -- {"$\bbQ(\zeta\sqrt[8]{2})$","$\bbQ(\zeta^3\sqrt[8]{2})$","$\bbQ(\zeta\sqrt[8]{2})$","$\bbQ(\zeta^2\sqrt[8]{2})$","$\bbQ(\sqrt[8]{2})$","$\bbQ(i,\sqrt[4]{2})$"};

        "$\bbQ(\zeta\sqrt[8]{2})$" -- "$\bbQ(i\sqrt[4]{2})$";

        "$\bbQ(\zeta^3\sqrt[8]{2})$" -- "$\bbQ(i\sqrt[4]{2})$";

        "$\bbQ(\zeta^2\sqrt[8]{2})$" -- "$\bbQ(\sqrt[4]{2})$";

        "$\bbQ(\sqrt[8]{2})$" -- "$\bbQ(\sqrt[4]{2})$";

        "$\bbQ(i,\sqrt[4]{2})$" -- {"$\bbQ(i\sqrt[4]{2})$","$\bbQ(\sqrt[4]{2})$","$\bbQ(i,\sqrt{2})$","$\bbQ((1+i)\sqrt[4]{2})$","$\bbQ((1-i)\sqrt[4]{2})$"};

        "$\bbQ(i\sqrt[4]{2})$" -- "$\bbQ(\sqrt{2})$";

        "$\bbQ(\sqrt[4]{2})$" -- "$\bbQ(\sqrt{2})$";

        "$\bbQ(i,\sqrt{2})$" -- {"$\bbQ(\sqrt{2})$","$\bbQ(i)$","$\bbQ(\sqrt{-2})$"};

        "$\bbQ((1+i)\sqrt[4]{2})$" -- "$\bbQ(\sqrt{-2})$";

        "$\bbQ((1-i)\sqrt[4]{2})$" -- "$\bbQ(\sqrt{-2})$";

        "$\bbQ(\sqrt{2})$" -- "$\bbQ$";

        "$\bbQ(i)$" -- "$\bbQ$";
        
        "$\bbQ(\sqrt{-2})$" -- "$\bbQ$";
      };
    \end{tikzpicture}
  \]  

  \[
    \begin{tikzpicture}
      % load the arrows library
      \usetikzlibrary{arrows.meta}
      % help lines (comment out when done) 
      \draw [help lines] (-1,-2) grid (5,6);
      % the triangle for the inclined plane
      \draw [thick] (0,0) -- (4,0) -- (4,3) -- (0,0);
      % the little right triangle corner 
      \draw (3.8,0) -- (3.8, .2) -- (4,.2);
      % the x measuring bar 
      \draw [{Bar}-{Bar}] (0,-0.3) -- (4,-0.3);
      % the y measuring bar
      \draw [{Bar}-{Bar}] (4.3,0) -- (4.3,3);
      % the little arc for the angle 
      \draw [radius = 0.5 cm] (0.5,0) arc
      [start angle=0, end angle = 36.87];
      \end{tikzpicture}
  \]

  \[
    \begin{tikzpicture} % load the arrows library
      \usetikzlibrary{arrows.meta}
      % help lines (comment out when done) 
      \draw [help lines] (-1,-2) grid (5,6);
      % the triangle for the inclined plane 
      \draw [thick] (0,0) -- (4,0) -- (4,3) -- (0,0);
      % the little right triangle corner
      \draw (3.8,0) -- (3.8, .2) -- (4,.2);
      % the x measuring bar 
      \draw [{Bar}-{Bar}] (0,-0.3) -- (4,-0.3);
      % the y measuring bar 
      \draw [{Bar}-{Bar}] (4.3,0) -- (4.3,3);
      % the little arc for the angle \draw [radius = 0.5 cm] (0.5,0) arc
      [start angle=0, end angle = 36.87];
      % the dashed line pointing about 5:00 from the apple 
      \draw [dashed, ->] (5/2,15/8) -- (79/25,199/200);
      % the line pointing down from the apple 
      \draw [->] (2.5,15/8) -- (5/2,.5);
% the line parallel to the plane 
\draw [dashed, ->] (79/25,199/200) -- (5/2,1/2);
\end{tikzpicture}
  \]

  \[
\begin{tikzpicture} % the dashed line pointing about 5:00 from the apple 
  \draw [dashed, -{Triangle}] (5/2,15/8) -- (79/25,199/200);
% the line pointing down from the apple
\draw [-{Triangle}] (2.5,15/8) -- (5/2,.5);
% the line parallel to the plane 
\draw [dashed, -{Triangle}] (79/25,199/200) -- (5/2,1/2); \end{tikzpicture} \]

\[ \begin{tikzpicture} % the dashed line pointing about 5:00 from the apple
  \draw [dashed, -{Triangle}] (5/2,15/8) -- (79/25,199/200);
% the line pointing down from the apple 
\draw [-{Triangle}] (2.5,15/8) -- (5/2,.5+.05);
% the line parallel to the plane 
\draw [dashed, -{Triangle}] (79/25,199/200) -- (5/2,1/2);
\end{tikzpicture}
\]

\[ \begin{tikzpicture}
% the dashed line pointing about 5:00 from the apple 
\draw [dashed, -{Stealth}] (5/2,15/8) -- (79/25,199/200);
% the line pointing down from the apple 
\draw [-{Stealth}] (2.5,15/8) -- (5/2,.5);
% the line parallel to the plane
\draw [dashed, -{Stealth}] (79/25,199/200) -- (5/2,1/2);
\end{tikzpicture} \]

\[
  \begin{tikzpicture}
    % load the arrows library \usetikzlibrary{arrows.meta}
    % help lines (comment out when done) 
    \draw [help lines] (-1,-2) grid (5,6);
    % the triangle for the inclined plane
    \draw [thick] (0,0) -- (4,0) -- (4,3) -- (0,0);
% the little right triangle corner 
\draw (3.8,0) -- (3.8, .2) -- (4,.2);
% the x measuring bar 
\draw [{Bar}-{Bar}] (0,-0.3) -- (4,-0.3);
% the y measuring bar
\draw [{Bar}-{Bar}] (4.3,0) -- (4.3,3);
% the little arc for the angle 
\draw [radius = 0.5 cm] (0.5,0) arc
[start angle=0, end angle = 36.87];
% the dashed line pointing about 5:00 from the apple, and the label on it
\draw [dashed, -{Stealth}] (5/2,15/8) -- (79/25,199/200); 
\node [above right=-2pt] at (283/100, 287/200) {$F_{\perp}$};
% the line pointing down from the apple, and the label on it
\draw [-{Stealth}] (2.5,15/8) -- (5/2,.5); 
\node [left] at (5/2, 19/16) {$F_g$};
% the line parallel to the plane, and the label on it 
\draw [dashed, -{Stealth}] (79/25,199/200) -- (5/2,1/2);
\node [below right=-1pt] at (283/100, 299/400) {$F_{||}$};
\end{tikzpicture}
\]

\begin{tikzpicture}
  \draw (0,0) circle (1ex);
\end{tikzpicture}

圆形
\tikz \draw (0,0) circle (1ex);

\tikz \fill (0,0) circle (1ex);

\tikz \filldraw[thick,fill=gray] (0,0) circle (0.5cm);

\begin{tikzpicture}
  \draw[->] (-1.2,0) -- (1.2,0);
  \draw[->] (0,-1.2) -- (0,1.2);
  \draw[thick] (0,0) circle (1);
\end{tikzpicture}

\tikz\draw (30:0.5) arc (30:150:0.5);

\begin{tikzpicture}
  \draw[thick] (0,0) circle (1);
  \fill[fill=gray,fill opacity=0.3] (0,0) -- (0:1) arc (0:120:1) -- cycle;
  \filldraw[fill=gray,fill opacity=0.5] (0,0) -- (0:0.3) arc (0:120:0.3) -- cycle;
\end{tikzpicture}

\begin{tikzpicture}
  \draw (0,0) circle (1);
  \fill (120:1) circle (2pt);
  \node[above left] (P) at (120:1) {$P$};
\end{tikzpicture}

\begin{tikzpicture}
  \draw (0,0) circle (1);
  \coordinate[label=120:$P$] (P) at (120:1);
  \draw (0,0) -- (P);
\end{tikzpicture}

\begin{tikzpicture}
  \newcommand\iangle{120}
  \draw[->] (-1.2,0) -- (1.2,0);
  \draw[->] (0,-1.2) -- (0,1.2);
  \draw[thick] (0,0) circle (1);
  \coordinate[label=\iangle:$P$] (P) at (\iangle:1);
  \coordinate[label=below:$P_0$] (P0) at (P |- 0,0);
  \draw (0,0) -- (P);
  \draw (P) -- (P0);
  \fill[fill=gray,fill opacity=0.2]
    (0,0) -- (0:1) arc (0:\iangle:1) -- cycle;
  \filldraw[fill=gray,fill opacity=0.5]
    (0,0) -- (0:0.3) arc (0:\iangle:0.3) -- cycle;
  \node[right] at (\iangle/2:0.3) {\ang{\iangle}};
\end{tikzpicture}

\begin{tikzpicture}
  \draw[domain=0:rad(360)] plot (\x,{sin(\x r)});
\end{tikzpicture}

\begin{tikzpicture}
  \draw (0,0) -- (2,0) node[right] {右};
  \draw (0,-1) -- node[above] {连线} (2,-1);
\end{tikzpicture}

\begin{tikzpicture}
  \draw[->] (0,0) -- ({rad(210)},0);
  \draw[->] (0,-1.2) -- (0,1.2);

  \foreach \t in {0,90,180} {
    \draw ({rad(\t)},-0.05) -- ({rad(\t)},0.05);
    \node[below,outer sep=2pt,fill=white,font=\small] at ({rad(\t)},0) {\ang{\t}};
  }
  \foreach \y in {-1,1} {\draw (-0.05,\y) -- (0.05,\y);}
\end{tikzpicture}

\begin{tikzpicture}[fill=lightgray]
  \begin{scope}[thick,->,fill=gray,xshift=-3cm]
    \filldraw (0,0) circle (1);
    \draw (-1.2,0) -- (1.2,0);
    \draw (0,-1.2) -- (0,1.2);
  \end{scope}
  \filldraw (0,0) circle (1);
\end{tikzpicture}

\begin{figure}
  \centering
  % \usepackage{siunitx}
  \begin{tikzpicture}[scale=1.5] % 整体放大坐标，但不影响字号
    \newcommand\iangle{120}
    % 左边的单位圆
    \begin{scope}[xshift=-2cm]
      \draw[->] (-1.2,0) -- (1.2,0);
      \draw[->] (0,-1.2) -- (0,1.2);
      \draw[thick] (0,0) circle (1);
      \coordinate[label=\iangle:$P$] (P) at (\iangle:1);
      \coordinate[label=below:$P_0$] (P0) at (P |- 0,0);
      \draw (0,0) -- (P);
      \draw (P) -- (P0);
      \fill[fill=gray,fill opacity=0.2]
        (0,0) -- (0:1) arc (0:\iangle:1) -- cycle;
      \filldraw[fill=gray,fill opacity=0.5]
        (0,0) -- (0:0.3) arc (0:\iangle:0.3) -- cycle;
      \node[right] at (\iangle/2:0.3) {\ang{\iangle}};
    \end{scope}
    % 右边的函数图
    \draw[->] (0,0) -- ({rad(210)}, 0);
    \draw[->] (0,-1.2) -- (0,1.2);
    \draw[thick, domain=0:rad(210)] plot (\x, {sin(\x r)})
      node[right] {$\sin x$};
    \foreach \t in {0, 90, 180} {
      \draw ({rad(\t)}, -0.05) -- ({rad(\t)}, 0.05);
      \node[below, outer sep=2pt, fill=white, font=\small]
        at ({rad(\t)}, 0) {\ang{\t}};
    }
    \foreach \y in {-1,1} {\draw (-0.05,\y) -- (0.05,\y);}
    \coordinate[label=above:$Q$] (Q) at ({rad(\iangle)}, {sin(\iangle)});
    \coordinate[label=below:$Q_0$] (Q0) at (Q |- 0,0);
    \draw (Q) -- (Q0);
    \draw[dashed] (P) -- (Q);
  \end{tikzpicture}
  \caption{正弦函数与单位圆（\textsf{TikZ} 实现）}
  \label{fig:tikzsine}
  \end{figure}
\end{document}